% !TEX root = main.tex
\clearpage
\section{Code guide: BachattTokenizer v0}

The tokenizer workstream contains three programs:

\renewcommand{\arraystretch}{1.2}
\begin{tabular}{@{}p{0.35\textwidth}p{0.51\textwidth}@{}}
  \textbf{File} & \textbf{Responsibility} \\
  \hline
  \texttt{train\_tokenizer.py} & Define and train the byte-level
  BPE tokenizer. \\
  \texttt{sample\_corpus.py} & Build a balanced text sample from cleaned
  domain pools for tokenizer training. \\
  \texttt{eval\_fertility.py} & Evaluate tokenizer efficiency/fertility
  by language. \\
\end{tabular}

Also, in parallel:
\renewcommand{\arraystretch}{1.2}

\subsection{\texttt{train\_tokenizer.py}}

\begin{lstlisting}[style=shadedpython,language={},basicstyle=\scriptsize\ttfamily,breaklines=true]
python train_tokenizer.py
    --input "data/*.txt"
    --vocab-size 160000
    --out tokenizer.json

          |
          v

ARGPARSE
Reads command-line settings:

--input       -> which text files to train on
--vocab-size  -> desired vocabulary size
--out         -> where to save tokenizer

          |
          v

GLOB FINDS FILES
glob.glob(args.input)

Example:
"data/*.txt"

          |
          v

[
  "data/english.txt",
  "data/hindi.txt",
  "data/tamil.txt"
]

These files become the tokenizer's
training corpus.

          |
          v

build_tokenizer()

\end{lstlisting}

\subsubsection*{Creates the tokenizer architecture}

\begin{lstlisting}[style=shadedpython,language={},basicstyle=\scriptsize\ttfamily,breaklines=true,escapeinside={(*}{*)}]
BPE TOKENIZER

Tokenizer(
    models.BPE(...)
)

BPE will learn frequently occurring
text pieces and merge them into tokens.

              |
              v

PRE-TOKENIZATION PIPELINE (IN ORDER):

+-------------------------------+
| 1. DIGITS                     |
|                               |
| individual_digits=True        |
|                               |
| 2026 -> 2 | 0 | 2 | 6         |
|                               |
| Makes number representation   |
| more consistent.              |
+---------------+---------------+
                |
                v

+-------------------------------+
| 2. WORD_PATTERN               |
|                               |
| Splits text into sensible     |
| BPE regions:                  |
|                               |
| * words                       |
| * punctuation                 |
| * whitespace                  |
| * letters + matras together   |
|                               |
| Example:                      |
|                               |
| "Hello, (*\dev{दुनिया}*)!"               |
|       |                       |
|       v                       |
| "Hello" | "," | " (*\dev{दुनिया}*)" | "!" |
|                               |
| Important: creates boundaries |
| that BPE cannot cross.        |
+---------------+---------------+
                |
                v

+-------------------------------+
| 3. BYTELEVEL                  |
|                               |
| Converts Unicode text into    |
| byte-level symbols.           |
|                               |
| English / Hindi / Tamil /     |
| emoji / etc. can all be       |
| represented.                  |
|                               |
| Helps avoid unknown chars.    |
+---------------+---------------+
                |
                v

BPE TRAINER

trainer = BpeTrainer(...)

Tells BPE HOW to learn vocabulary.

              |
              +-- vocab_size = 160000
              |   "Stop around this vocabulary size"
              |
              +-- min_frequency = 2
              |   "Ignore extremely rare merge candidates"
              |
              +-- special_tokens
              |   <|bos|>
              |   <|eos|>
              |   <|pad|>
              |
              +-- reserved tokens
              |   <|reserved_0|>
              |   ...
              |   <|reserved_63|>
              |
              +-- byte alphabet
                  Include basic byte-level
                  symbols from the start.
              |
              v
           tok.train(files, trainer)

\end{lstlisting}

\subsubsection*{Actual BPE training happens here}

\begin{lstlisting}[style=shadedpython,language={},basicstyle=\scriptsize\ttfamily,breaklines=true]
Training text
      |
      v
pre-tokenization
      |
      v
count frequently occurring adjacent pieces
      |
      v
merge frequent pairs
      |
      v
repeat many times
      |
      v
learn larger useful tokens

\end{lstlisting}

\subsubsection*{Simplified example: how BPE merges text into tokens}

\begin{lstlisting}[style=shadedpython,language={},basicstyle=\scriptsize\ttfamily,breaklines=true]
h | e | l | l | o
      |
      v
    l + l
      |
      v

h | e | ll | o
      |
      v
learn more merges
      |
      v
    "hello"
      |
      v
BUILD VOCABULARY

\end{lstlisting}

\subsubsection*{Vocabulary: learned tokens, special tokens, reserved
tokens, and byte alphabet}

\begin{lstlisting}[style=shadedpython,language={},basicstyle=\scriptsize\ttfamily,breaklines=true,escapeinside={(*}{*)}]

SPECIAL TOKENS
-> fixed control tokens

<|bos|>
<|eos|>
<|pad|>

RESERVED SLOTS
-> saved for future features

<|reserved_0|>
...
<|reserved_63|>

BYTE ALPHABET
-> guarantees basic byte-level coverage

LEARNED BPE TOKENS
-> common pieces learned from the data

Examples might include:

"the"
"ing"
"hello"
" (*\dev{भारत}*)"
etc.

    |
    v

tok.save(args.out)
Saves everything needed to reuse the tokenizer:

* vocabulary
* BPE merge rules
* pre-tokenizer setup
* decoder setup
* special tokens

      |
      v
tokenizer.json
(This is the finished trained tokenizer)

\end{lstlisting}

\subsubsection*{Later Usage: Encoding and Decoding}

\begin{lstlisting}[style=shadedpython,language={},basicstyle=\scriptsize\ttfamily,breaklines=true]

TEXT
  |
  v
tokenizer.json
  |
  v
TOKEN IDs

and:

TOKEN IDs
  |
  v
tokenizer.json
  |
  v
TEXT
\end{lstlisting}

\subsubsection*{What is BPE?}

\textbf{BPE} stands for \textbf{Byte Pair Encoding}. It is a data
compression technique that has become the foundational
\textbf{tokenization algorithm} used by modern Large Language Models
(LLMs) like GPT-4, Llama, and Claude to break down human text into
numbers the computer can understand.

Instead of splitting text into whole words (which creates a massive
dictionary) or individual letters (which takes too many steps to
read), BPE finds the perfect middle ground: \textbf{subwords}.

\subsubsection*{Word vs. Character vs. BPE Tokenization}

\begin{center}
  \renewcommand{\arraystretch}{1.4}
  \begin{tabular}{@{}p{0.18\textwidth}p{0.32\textwidth}p{0.22\textwidth}p{0.18\textwidth}@{}}
    \textbf{Tokenizer Type} & \textbf{How it splits ``Unbelievable''}
    & \textbf{Pros} & \textbf{Cons} \\
    \hline
    \textbf{Word-Based} & \texttt{["Unbelievable"]} & Easy to
    understand. & Cannot handle typos or new words (OOV error). \\
    \textbf{Character-Based} & \texttt{["U", "n", "b", "e", "l",
    ...]} & Small vocabulary, can spell anything. & Sequences become
    too long; loses semantic meaning per token. \\
    \textbf{BPE (Subword)} & \texttt{["Un", "believ", "able"]} &
    Handles rare words by breaking into familiar roots; highly
    efficient. & Text is split in ways that don't always align with
    human grammar. \\
  \end{tabular}
\end{center}

\subsubsection*{Why is BPE crucial for LLMs?}

\begin{itemize}
    \setlength{\itemsep}{0pt}
  \item \textbf{It Solves Unknown Words:} If an AI encounters a brand
    new word like \texttt{AI-ification}, it won't crash. It will
    break it down into pieces it \textit{does} know: \texttt{["AI",
    "-", "ific", "ation"]}.
  \item \textbf{It is Multilingual:} BPE can ingest bytes or
    characters from any language, dynamically building tokens for
    Hindi, Chinese, or English based on frequency.
  \item \textbf{It Saves Memory:} By grouping common sequences like
    \texttt{" the"} or \texttt{" ing"} into single tokens, the model
    can process much longer essays or code files within its limited
    context window.
\end{itemize}

\subsubsection*{The training script}

The whole script does this:
\begin{lstlisting}[style=shadedpython,language={},basicstyle=\footnotesize\ttfamily]
TEXT FILES
    |
    v
find files matching --input
    |
    v
build tokenizer
    |
    v
pre-tokenize text
    |--- split digits individually
    |--- split words / punctuation / whitespace
    |--- convert characters to byte-level symbols
    |
    v
train BPE merges
    |
    v
build vocabulary
    |--- normal learned tokens
    |--- special tokens
    |--- reserved tokens
    |
    v
save tokenizer.json
\end{lstlisting}

To run the script:
\begin{lstlisting}[style=shadedpython,breaklines=true]
python train_tokenizer.py \
  --input "path/to/*.txt" \
  --vocab-size 160000 \
  --out tokenizer-160k.json
\end{lstlisting}

\subsubsection*{Training configuration}

\texttt{build\_tokenizer()} assembles three pre-tokenizers in sequence:

\begin{lstlisting}[style=shadedpython,breaklines=true]
  # each digit is its own token: 2026 -> 2 0 2 6
  pre_tokenizers.Digits(individual_digits=True),

  # creates linguistically useful word-like segments for BPE training
  pre_tokenizers.Split(WORD_PATTERN, behavior="isolated"),

  # all remaining characters to byte-level symbols, to avoid unknown-token failure
  # set add_prefix_space=False to avoid automatically space insert at the beginning
  # set use_regex=False to avoid ByteLevel perform its own regex splitting
  pre_tokenizers.ByteLevel(add_prefix_space=False, use_regex=False),
\end{lstlisting}

\noindent\textbf{Special tokens}
\begin{lstlisting}[style=shadedpython,breaklines=true]
SPECIAL_TOKENS = ["<|bos|>", "<|eos|>", "<|pad|>"]
\end{lstlisting}
These are tokens with special meanings:

\begin{lstlisting}[style=shadedpython,language={},basicstyle=\footnotesize\ttfamily]
<|bos|> --- beginning of sequence
<|eos|> --- end of sequence
<|pad|> --- padding
\end{lstlisting}

For example: \texttt{< |bos| > Hello world < |eos| >}. They are not normal
words learned by BPE. They are inserted directly into the vocabulary.

\noindent\textbf{Reserved tokens}
\begin{lstlisting}[style=shadedpython,breaklines=true]
RESERVED = [f"<|reserved_{i}|>" for i in range(64)]
\end{lstlisting}
This creates:
\begin{lstlisting}[style=shadedpython,language={},basicstyle=\footnotesize\ttfamily]
<|reserved_0|>
<|reserved_1|>
...
<|reserved_63|>
\end{lstlisting}

Why? Suppose today you have:
\begin{lstlisting}[style=shadedpython,language={},basicstyle=\footnotesize\ttfamily]
ID 0 → <|bos|>
ID 1 → <|eos|>
ID 2 → <|pad|>
\end{lstlisting}
Later you decide to add \texttt{<|tool\_call|>}, \texttt{<|fim\_start|>},
\texttt{<|vision|>}. If you insert new tokens later, vocabulary IDs might
shift. That can break compatibility with previously trained models. So
instead, you reserve slots now.

\subsubsection*{The matra-aware split}

The central tokenizer finding is encoded in \texttt{WORD\_PATTERN}:

\begin{lstlisting}[style=shadedpython,breaklines=true]
WORD_PATTERN = Regex(
    r" ?[\p{L}\p{M}]+| ?[^\s\p{L}\p{M}\p{N}]+|\s+"
)
\end{lstlisting}

Think of the regex as a \textbf{3-way text splitter}:

\begin{lstlisting}[style=shadedpython,language={},basicstyle=\footnotesize\ttfamily]
WORD_PATTERN
    |
    v
Input text
    |
    +----------------------+----------------------+------------------+
    |                      |                      |
    v                      v                      v
  RULE 1                 RULE 2                 RULE 3
  WORDS                  SYMBOLS               WHITESPACE

 ?[\p{L}\p{M}]+      ?[^\s\p{L}\p{M}\p{N}]+      \s+
    |                      |                      |
    v                      v                      v
optional space         optional space          one or more
+ letters              + punctuation           whitespace chars
                       / symbols

Examples:
" hello"               " !!!"                  "\n\n"
"world"                 ","                     "\t"
" GPT"                  "??"                    "   "
\end{lstlisting}

\noindent\textbf{1.\ }\verb| ?[\p{L}\p{M}]+|

This matches a word made of letters, optionally including one space before
it. Break it down:

\begin{itemize}
    \setlength{\itemsep}{0pt}
    \item \verb| ?| --- one space, optionally. The \texttt{?} means zero or one
    occurrence.
    \item \verb|[\p{L}\p{M}]| --- \verb|\p{L}| is a Unicode letter;
    \verb|\p{M}| is a Unicode combining mark.
  \item \texttt{+} --- one or more.
\end{itemize}

So \verb| ?[\p{L}\p{M}]+| means roughly: optional leading space + one or
more letters/marks. This forces to include the matra (vowel sign) with its
base letter, which is crucial for Indic scripts.

Example: \texttt{"hello world"} could match
\begin{lstlisting}[style=shadedpython,breaklines=true]
  "hello" and " world"
\end{lstlisting}

Notice that the space before
\texttt{"world"} becomes part of the token.

The key Unicode categories are:
\begin{itemize}
    \setlength{\itemsep}{0pt}
    \item \verb|\p{L}| = Letter
    \item \verb|\p{M}| = Mark
\end{itemize}
A matra is usually stored as a combining mark, so it falls under
\verb|\p{M}|. For example, conceptually:
\begin{center}
  \renewcommand{\arraystretch}{1.3}
  \begin{tabular}{@{}lll@{}}
    \dev{क} & = base consonant &
    $\rightarrow$ \verb|\p{L}| \\
    \dev{ा} & = matra & $\rightarrow$ \verb|\p{M}| \\
  \end{tabular}
\end{center}
Together they visually form: \dev{का}.

This is important because BPE is usually only allowed to merge symbols
inside the segment it receives. So if preprocessing gives BPE:
\begin{center}
  \begin{tabular}{@{}ll@{}}
    segment 1: & ``\dev{क}'' \\
    segment 2: & ``\dev{ा}'' \\
  \end{tabular}
\end{center}
BPE cannot later say ``Actually, let me merge
\dev{क}~+~\dev{ा}
into \dev{का}'' because the boundary has already
been imposed.

\noindent\textbf{2.\ }\verb| ?[^\s\p{L}\p{M}\p{N}]+|

This matches punctuation and symbols, again optionally with one leading
space.

Start with \verb| ?| --- same meaning: optionally include one space
before the token.

Then \verb|[^ ... ]|: the \texttt{\textasciicircum} inside
square brackets means \emph{NOT} any of these characters:

\begin{itemize}
    \setlength{\itemsep}{0pt}
    \item \verb|\s| --- whitespace
    \item \verb|\p{L}| --- letters
    \item \verb|\p{M}| --- combining marks
    \item \verb|\p{N}| --- numbers
\end{itemize}

So \verb|[^\s\p{L}\p{M}\p{N}]| means anything that is \emph{not} whitespace, a
letter, a combining mark, or a number.

In practice, that means things like
\texttt{! ? . , : ; @ \# \$ ( )}. And \texttt{+} means one or more.

Example: \texttt{"hello!!!"} might become
\begin{lstlisting}[style=shadedpython,breaklines=true]
  "hello"
  "!!!"
\end{lstlisting}
Example: \texttt{"hello !!!"} might become
\begin{lstlisting}[style=shadedpython,breaklines=true]
  "hello"
  " !!!"
\end{lstlisting}
because the optional space can attach to the punctuation token.

\noindent\textbf{3.\ }\verb|\s+|

This matches one or more whitespace characters: \verb|\s| is whitespace and
\texttt{+} is one or more. Whitespace includes space, tab, and newline. So
\texttt{"\textbackslash n\textbackslash n"} can be matched as one whitespace
chunk.

\subsection{\texttt{sample\_corpus.py}}

\begin{lstlisting}[style=shadedpython,language={},basicstyle=\scriptsize\ttfamily,breaklines=true]
02-data-pipeline/pools/
|-- global_english/     raw.jsonl + clean.jsonl      <- FineWeb-Edu
|-- indian_english/     sangraha_eng, wiki_india,    <- Sangraha verified/eng + en-wiki India topics
|                       raw.jsonl + clean.jsonl
|-- indian_languages/   sangraha_hin.jsonl,          <- Sangraha verified, 11 languages
|                       sangraha_ben.jsonl, _tam,       (hin ben tam tel mar guj kan mal pan ori asm)
|                       ... + raw.jsonl + clean.jsonl
|-- code/               raw.jsonl + clean.jsonl      <- 4 allowlisted GitHub repos (.py files)
|-- finance_econ_law/   wiki_en, wiki_hi,            <- curated Wikipedia titles (75 en + 25 hi)
|                       raw.jsonl + clean.jsonl
|-- math_reasoning/     raw.jsonl + clean.jsonl      <- curated Wikipedia titles (40 en)
\end{lstlisting}

How each pool was fetched --- \texttt{ingest/run\_ingest.sh} drives three
fetchers:
\begin{enumerate}
    \setlength{\itemsep}{0pt}
  \item \texttt{ingest\_hf\_rows.py} --- pages through the HuggingFace
    datasets-server rows API (100 rows/request, polite pacing, 429
    backoff). Pulled 3{,}000 FineWeb-Edu docs, $\sim$600 Sangraha
    English, and $\sim$300 per Indic language.
  \item \texttt{ingest\_wikipedia.py} --- MediaWiki extracts API against
    the curated title lists in \texttt{ingest/titles\_*.txt}
    (finance, math, India topics; one full-text extract per request
    --- the API silently caps batches at one).
  \item \texttt{ingest\_github\_code.py} --- codeload tarballs of four
    permissive repos (requests/flask/click/more-itertools), one doc
    per \texttt{.py} file, license recorded per doc.
\end{enumerate}
Every source has a row in \texttt{legal\_register.md} before ingest
(S010--S030).

\subsubsection{Data Ingestion}

\begin{paragraphblock}{Create the six data pools}
  The script creates:
\begin{lstlisting}[style=shadedpython,language={},basicstyle=\scriptsize\ttfamily]
pools/
|-- global_english/
|-- indian_english/
|-- indian_languages/
|-- code/
|-- finance_econ_law/
|-- math_reasoning/
\end{lstlisting}
  Each pool has a different purpose.

  \begin{center}
    \renewcommand{\arraystretch}{1.2}
    \small
    \begin{tabular}{@{}l r r r@{}}
      \textbf{Pool} & \textbf{Charter range} & \textbf{Working weight} &
      \textbf{At 7T (Beta)} \\
      \hline
      \texttt{global\_english} & 45--55\% & 50.0\% & 3.5T \\
      \texttt{indian\_english} & 15--20\% & 17.5\% & 1.23T \\
      \texttt{indian\_languages} & 10--15\% & 12.5\% & 0.88T \\
      \texttt{code} & 10--15\% & 10.0\% & 0.70T \\
      \texttt{finance\_econ\_law} & 5--10\% & 7.5\% & 0.53T \\
      \texttt{math\_reasoning} & 2--5\% & 2.5\% & 0.18T \\
      \hline
      & & 100\% & 7T \\
    \end{tabular}
  \end{center}

  Notes:
  \begin{itemize}
      \setlength{\itemsep}{0pt}
    \item Within \texttt{indian\_languages} there's no per-language
      weighting yet --- currently $\sim$equal pulls ($\sim$300 docs
      each of the 11 languages). A real policy (e.g., weight by speaker
      population vs equal-per-language vs temperature sampling) is an
      open decision; it matters more for the tokenizer than the model
      right now, which is why \texttt{sample\_corpus.py} balanced
      per-language for tokenizer training.
  \end{itemize}
\end{paragraphblock}

\begin{paragraphblock}{Global English}
  Source: \texttt{HuggingFaceFW/fineweb-edu}, config \texttt{sample-10BT}.

\begin{lstlisting}[style=shadedpython,language={},basicstyle=\scriptsize\ttfamily]
FineWeb-Edu
    |
    v
take ~3000 rows by default
    |
    v
pools/global_english/raw.jsonl
\end{lstlisting}

  Code:
\begin{lstlisting}[style=shadedpython,breaklines=true]
ingest_hf_rows.py \
    --dataset HuggingFaceFW/fineweb-edu \
    --config sample-10BT \
    --split train \
    --rows "$ROWS_FINEWEB"
\end{lstlisting}
  Purpose: provide broad, high-quality general English text.
\end{paragraphblock}

\begin{paragraphblock}{Indian English}
  This combines two sources: Sangraha verified English + English Wikipedia
  pages about India.

  Flow:
\begin{lstlisting}[style=shadedpython,language={},basicstyle=\scriptsize\ttfamily]
ai4bharat/sangraha
verified / eng
        |
        v
sangraha_eng.jsonl

             +

curated English Wikipedia
India-related titles
        |
        v
wiki_india.jsonl

             | concatenate
             v

pools/indian_english/raw.jsonl
\end{lstlisting}
  So Indian English gets both general India-focused English from Sangraha
  and curated India-specific Wikipedia content.
\end{paragraphblock}

\begin{paragraphblock}{Indian languages}
  Source: AI4Bharat Sangraha, config \texttt{verified}.

  Languages:
  \begin{center}
    \renewcommand{\arraystretch}{1.15}
    \small
    \begin{tabular}{@{}ll@{}}
      \texttt{hin} $\rightarrow$ Hindi &
      \texttt{ben} $\rightarrow$ Bengali \\
      \texttt{tam} $\rightarrow$ Tamil &
      \texttt{tel} $\rightarrow$ Telugu \\
      \texttt{mar} $\rightarrow$ Marathi &
      \texttt{guj} $\rightarrow$ Gujarati \\
      \texttt{kan} $\rightarrow$ Kannada &
      \texttt{mal} $\rightarrow$ Malayalam \\
      \texttt{pan} $\rightarrow$ Punjabi &
      \texttt{ori} $\rightarrow$ Odia \\
      \texttt{asm} $\rightarrow$ Assamese &
    \end{tabular}
  \end{center}

  The script loops through all 11:
\begin{lstlisting}[style=shadedpython,language={},basicstyle=\scriptsize\ttfamily]
Hindi       -> sangraha_hin.jsonl
Bengali     -> sangraha_ben.jsonl
Tamil       -> sangraha_tam.jsonl
...
Assamese    -> sangraha_asm.jsonl
\end{lstlisting}
  Then concatenates all of them:
\begin{lstlisting}[style=shadedpython,language={},basicstyle=\scriptsize\ttfamily]
11 language files
        |
        v
pools/indian_languages/raw.jsonl
\end{lstlisting}
  By default it takes 300 rows per language unless overridden.
\end{paragraphblock}

\begin{paragraphblock}{Code}
  Source: allowlisted permissive GitHub repositories.

  The script calls \texttt{ingest/ingest\_github\_code.py} and writes
  \texttt{pools/code/raw.jsonl}.

  The important idea: code comes only from an allowlisted set of
  repositories intended to satisfy the project's licensing rules. The
  exact repo list is not shown in this script.
\end{paragraphblock}

\begin{paragraphblock}{Finance / economics / law}
  This is built from curated Wikipedia titles in English and Hindi.

  Flow:
\begin{lstlisting}[style=shadedpython,language={},basicstyle=\scriptsize\ttfamily]
titles_finance_en.txt
        |
        v
Wikipedia EN
        |
        v
wiki_en.jsonl

titles_finance_hi.txt
        |
        v
Wikipedia HI
        |
        v
wiki_hi.jsonl

        | concatenate
        v

pools/finance_econ_law/raw.jsonl
\end{lstlisting}
  So the domain content is deliberately selected rather than pulling
  arbitrary Wikipedia pages.
\end{paragraphblock}

\begin{paragraphblock}{Math / reasoning}
  Source: curated English Wikipedia pages. The title list comes from
  \texttt{ingest/titles\_math\_en.txt}.

  Flow:
\begin{lstlisting}[style=shadedpython,language={},basicstyle=\scriptsize\ttfamily]
curated math titles
       |
       v
English Wikipedia
       |
       v
pools/math_reasoning/raw.jsonl
\end{lstlisting}
\end{paragraphblock}

How \texttt{raw.jsonl} became \texttt{clean.jsonl} --- each pool ran
through the four-stage pipe:
\begin{lstlisting}[style=shadedpython,language={},basicstyle=\scriptsize\ttfamily,breaklines=true]
raw.jsonl -> langid.py -> filters.py -> dedup_minhash.py -> decontaminate.py -> clean.jsonl
            lang/script/  Gopher-style  MinHash+LSH        13-gram overlap
            code-mix tag  quality rules near-dup removal   vs BachattBench
\end{lstlisting}

\subsubsection{Langid}

\texttt{langid} adds language, script, and code-mix metadata to every
document.

Input document:
\begin{lstlisting}[style=shadedpython,breaklines=true,escapeinside={(*}{*)}]
{
  "id": "...",
  "source": "S011",
  "text": "(*\dev{भारत}*) ki economy badh rahi hai"
}
\end{lstlisting}
The script adds:
\begin{lstlisting}[style=shadedpython,breaklines=true]
{
  "lang": "...",
  "script": "...",
  "code_mix": true
}
\end{lstlisting}
So the document becomes richer.

\begin{paragraphblock}{script\_profile(text)}
  This function goes character by character. For every alphabetic character:
\begin{lstlisting}[style=shadedpython,language={},basicstyle=\scriptsize\ttfamily]
character
   |
   v
Unicode code point
   |
   v
which script range?
\end{lstlisting}
  The script knows ranges for Devanagari, Bengali, Gurmukhi, Gujarati, Odia,
  Tamil, Telugu, Kannada, Malayalam, and Latin.

  Example: \texttt{"}\dev{भारत}\texttt{ test"} might roughly produce
  \texttt{Devanagari → 0.60}, \texttt{Latin → 0.40}. The result is a
  dictionary of fractions:
\begin{lstlisting}[style=shadedpython,breaklines=true]
{
    "devanagari": 0.60,
    "latin": 0.40
}
\end{lstlisting}
\end{paragraphblock}

\begin{paragraphblock}{How native-script language is inferred}
  The largest script proportion wins: \texttt{top = max(prof, key=prof.get)}.
  If the top script is not Latin, this mapping is used:
  \begin{center}
    \renewcommand{\arraystretch}{1.15}
    \small
    \begin{tabular}{@{}ll@{}}
      Devanagari $\rightarrow$ Hindi &
      Bengali $\rightarrow$ Bengali \\
      Gurmukhi $\rightarrow$ Punjabi &
      Gujarati $\rightarrow$ Gujarati \\
      Odia $\rightarrow$ Odia &
      Tamil $\rightarrow$ Tamil \\
      Telugu $\rightarrow$ Telugu &
      Kannada $\rightarrow$ Kannada \\
      Malayalam $\rightarrow$ Malayalam &
    \end{tabular}
  \end{center}
  This is explicitly coarse.
\end{paragraphblock}

\begin{paragraphblock}{Code-mix detection for native-script text}
  If a document is mainly native script but contains substantial Latin text
  (\texttt{prof.get("latin", 0.0) > 0.15}), then more than 15\% Latin letters
  $\Rightarrow$ \texttt{code\_mix = True}.

  Example: \texttt{"}\dev{भारत}\texttt{ ki economy bahut fast grow kar rahi
  hai"}. If Devanagari is still dominant but Latin exceeds 15\%, it becomes:
\begin{lstlisting}[style=shadedpython,breaklines=true]
{
  "lang": "hindi",
  "script": "devanagari",
  "code_mix": true
}
\end{lstlisting}
\end{paragraphblock}

\begin{paragraphblock}{Latin text is harder}
  Example: \texttt{"What are you doing?"} vs \texttt{"Tum kya kar rahe ho?"}
  --- both use Latin letters.

  It defines:
\begin{lstlisting}[style=shadedpython,breaklines=true]
HINGLISH_MARKERS = {
    "hai", "hain", "nahi", "kya", "kaise",
    "karna", "chahiye", "mera", "mujhe",
    "aur", "lekin", ...
}
\end{lstlisting}
  Then:
  \[
    \texttt{marker\_frac}
    =\frac{\text{number of Hinglish marker words}}{\text{total words}}.
  \]
  If marker fraction $> 5\%$, it labels the document:
\begin{lstlisting}[style=shadedpython,breaklines=true]
{
  "lang": "hinglish",
  "script": "latin",
  "code_mix": true
}
\end{lstlisting}
  Otherwise:
\begin{lstlisting}[style=shadedpython,breaklines=true]
{
  "lang": "english",
  "script": "latin",
  "code_mix": false
}
\end{lstlisting}
\end{paragraphblock}

\subsubsection{Filters}

This is where documents can actually be dropped for quality. \\
\texttt{filters.py} applies eight independent tests and returns all failure
reasons rather than stopping at the first one:

\begin{center}
  \renewcommand{\arraystretch}{1.15}
  \small
  \begin{tabular}{@{}p{0.30\textwidth}p{0.55\textwidth}@{}}
    \textbf{Filter} & \textbf{Rejected condition} \\
    \hline
    Length & fewer than 50 or more than 200{,}000 words \\
    Mean word length & outside 2--12 characters \\
    Symbol ratio & too many hash signs or ellipses per word \\
    Bullet/ellipsis lines & page dominated by list or fragment patterns \\
    Duplicate lines & repeated lines exceed 30\% of characters \\
    Top bigram & one two-word phrase dominates the document \\
    Alphabetic words & fewer than 60\% of words contain a letter \\
    Stopwords & fewer than two known stopwords for English, Hindi, or Hinglish
  \end{tabular}
\end{center}

\begin{paragraphblock}{Length filter}
\begin{lstlisting}[style=shadedpython]
MIN_WORDS = 50
MAX_WORDS = 200_000
\end{lstlisting}
\end{paragraphblock}

\begin{paragraphblock}{Mean word length}
\begin{lstlisting}[style=shadedpython]
MIN_MEAN_WORD_LEN = 2.0
MAX_MEAN_WORD_LEN = 12.0
\end{lstlisting}
  It computes
  \[
    \frac{\text{sum of word lengths}}{\text{number of words}}.
  \]
  Example: \texttt{["this", "is", "good"]} has lengths $4+2+4=10$, so
  $10/3\approx 3.33$ --- passes. But something like \texttt{x x x x x x x}
  has suspiciously tiny average word length, while
  \texttt{asjdklasjdklasjdklasjdkl...} may have an absurdly high average.
\end{paragraphblock}

\begin{paragraphblock}{Symbol ratio}
\begin{lstlisting}[style=shadedpython]
MAX_SYMBOL_WORD_RATIO = 0.10
\end{lstlisting}
  It counts occurrences of \texttt{\#} and \texttt{...} and compares them to
  word count. If symbols are too frequent (\texttt{symbols / words > 10\%}),
  the document is rejected.
\end{paragraphblock}

\begin{paragraphblock}{Bullet / ellipsis line filter}
  It checks non-empty lines. If over 90\% start like \texttt{-}, \texttt{*},
  or \texttt{•}, then \texttt{bullet\_lines}. If over 30\% end in
  \texttt{...} or \texttt{…}, then \texttt{ellipsis\_lines}. This catches
  documents that look like
\begin{lstlisting}[style=shadedpython,language={},basicstyle=\scriptsize\ttfamily]
• item
• item
• item
• item
...
\end{lstlisting}
  rather than normal prose.
\end{paragraphblock}

\begin{paragraphblock}{Duplicate-line filter}
\begin{lstlisting}[style=shadedpython]
MAX_DUP_LINE_FRAC = 0.30
\end{lstlisting}
  It looks for repeated lines, but weights by character count. Why
  character-weighted? Because repeating a tiny heading
\begin{lstlisting}[style=shadedpython,language={},basicstyle=\scriptsize\ttfamily]
Chapter
Chapter
Chapter
\end{lstlisting}
  should matter less than repeating a giant paragraph. Conceptually:
  \[
    \frac{\text{repeated characters}}{\text{total line characters}}.
  \]
  If that exceeds 30\%: reject $\rightarrow$ \texttt{dup\_lines}.
\end{paragraphblock}

\begin{paragraphblock}{Repeated bigram filter}
  A bigram here means two consecutive words (e.g.\ \texttt{"the cat"}). The
  code counts every bigram. If one bigram occurs too often relative to the
  document:
\begin{lstlisting}[style=shadedpython]
max_bigram_count * 2 / total_words > 0.18
\end{lstlisting}
  then \texttt{repeated\_bigram}. This catches text like
  \texttt{buy now buy now buy now buy now buy now...} where one phrase is
  repeated unnaturally often.
\end{paragraphblock}

\begin{paragraphblock}{Alphabetic-word fraction}
\begin{lstlisting}[style=shadedpython]
MIN_ALPHA_WORD_FRAC = 0.60
\end{lstlisting}
  At least 60\% of whitespace-separated words must contain an alphabetic
  character. Example bad document: \texttt{12345 \$\$\$ --- 9987 \#\#\# 123}
  --- most ``words'' contain no letters, so it gets rejected.
\end{paragraphblock}

\begin{paragraphblock}{Stopword sanity check}
  For English: \texttt{the}, \texttt{and}, \texttt{of}, \texttt{to},
  \texttt{in}, \texttt{is}, \ldots\
  For Hindi: \dev{है}, \dev{और}, \dev{के}, \dev{में}, \dev{से}, \ldots\
  For Hinglish: \texttt{hai}, \texttt{aur}, \texttt{ke}, \texttt{mein},
  \ldots\

  The document must contain at least
\begin{lstlisting}[style=shadedpython]
MIN_STOPWORD_HITS = 2
\end{lstlisting}
  distinct stopwords from its claimed language list. This is a sanity check.
  Example: \texttt{lang = english} but the document contains no common
  English function words at all --- that may indicate a wrong language
  label, garbage, or non-prose content, so it gets rejected.
\end{paragraphblock}

\begin{paragraphblock}{\texttt{judge(doc)}}
  This line is elegant:
\begin{lstlisting}[style=shadedpython]
reasons = [r for f in FILTERS if (r := f(doc))]
\end{lstlisting}
  Meaning: run every filter, then collect every rejection reason. Example:
\begin{lstlisting}[style=shadedpython]
reasons = [
    "too_short:31",
    "few_stopwords"
]
\end{lstlisting}
  Then \texttt{return (not reasons, reasons)} --- so no reasons $\Rightarrow$
  \texttt{keep = True}; one or more reasons $\Rightarrow$
  \texttt{keep = False}.
\end{paragraphblock}

\subsubsection{Dedup\_minhash}

We are no longer asking ``Is this document good?'' We are asking ``Have we
already seen essentially the same document?''

The method is:
\begin{lstlisting}[style=shadedpython,language={},basicstyle=\scriptsize\ttfamily]
5-word shingles
    |
    v
MinHash signatures
    |
    v
LSH candidate search
    |
    v
exact Jaccard verification
    |
    v
drop near duplicates
\end{lstlisting}

\begin{paragraphblock}{Step 1: 5-word shingles}
\begin{lstlisting}[style=shadedpython]
SHINGLE = 5
\end{lstlisting}
  Suppose: \texttt{"the quick brown fox jumps over the lazy dog"}. 5-word
  shingles are:
\begin{lstlisting}[style=shadedpython,language={},basicstyle=\scriptsize\ttfamily]
the quick brown fox jumps
quick brown fox jumps over
brown fox jumps over the
fox jumps over the lazy
jumps over the lazy dog
\end{lstlisting}
  Each one gets hashed into an integer. So a document becomes a set:
\begin{lstlisting}[style=shadedpython]
{
 hash(shingle1),
 hash(shingle2),
 hash(shingle3),
 ...
}
\end{lstlisting}
  Why? Because overlapping 5-word sequences capture document similarity much
  better than exact whole-document matching.
\end{paragraphblock}

\begin{paragraphblock}{MinHash signatures}
  A document may contain thousands of shingles. Comparing every document
  against every other document would be expensive. So MinHash compresses the
  set into \texttt{NUM\_HASHES = 128} numbers. Conceptually:
\begin{lstlisting}[style=shadedpython,language={},basicstyle=\scriptsize\ttfamily]
thousands of shingles
        |
        v
128-value fingerprint
\end{lstlisting}
  Documents with similar shingle sets tend to have similar MinHash signatures.
\end{paragraphblock}

\begin{paragraphblock}{Fixed seed}
\begin{lstlisting}[style=shadedpython]
SEED = 20260814
\end{lstlisting}
  This makes generated hash-function coefficients reproducible. So the
  intention is: same data + same code + same seed $=$ same signatures.
\end{paragraphblock}

\begin{paragraphblock}{LSH banding}
  The 128-value signature is split into \texttt{BANDS = 16}. So each band
  contains $128 / 16 = 8$ values. Visually:
\begin{lstlisting}[style=shadedpython,language={},basicstyle=\scriptsize\ttfamily]
128-value signature

[8][8][8][8][8][8][8][8]
[8][8][8][8][8][8][8][8]

= 16 bands
\end{lstlisting}
  For each band, it creates a bucket key. If two documents share any identical
  band, they become candidate duplicates. This avoids comparing every pair.
\end{paragraphblock}

\begin{paragraphblock}{Why LSH?}
  Without LSH: 100{,}000 docs $\rightarrow$ potentially billions of
  comparisons. With LSH:
\begin{lstlisting}[style=shadedpython,language={},basicstyle=\scriptsize\ttfamily]
MinHash
  |
  v
bucket similar-looking docs together
  |
  v
only verify likely matches
\end{lstlisting}
  So MinHash approximates similarity; LSH efficiently finds likely similar
  pairs.
\end{paragraphblock}

\begin{paragraphblock}{Exact Jaccard verification}
  Candidate pairs are verified using \texttt{jaccard(a, b)}:
  \[
    J(A,B)=\frac{|A \cap B|}{|A \cup B|}.
  \]
  Example: $A=\{1,2,3,4,5\}$, $B=\{1,2,3,4,6\}$ gives
  $|A\cap B|=4$, $|A\cup B|=6$, so $J=4/6=0.667$.

  Threshold:
\begin{lstlisting}[style=shadedpython]
JACCARD_THRESHOLD = 0.8
\end{lstlisting}
  So Jaccard $\ge 0.8$ $\rightarrow$ near duplicate $\rightarrow$ drop later
  doc.
\end{paragraphblock}

\begin{paragraphblock}{First-seen wins}
  The implementation keeps the first document seen and rejects later
  near-duplicates:
\begin{lstlisting}[style=shadedpython,language={},basicstyle=\scriptsize\ttfamily]
Doc A
  | kept
  v
Doc B ~ A
  | dropped
  v
Doc C ~ A
  | dropped
\end{lstlisting}
  So file order matters here too.
\end{paragraphblock}

\subsubsection{Decontaminate}

This stage has a different goal: \textbf{prevent benchmark leakage}.

It creates a set of sequences from the evaluation data and removes any
training document sharing one of them. The configured size is
\begin{lstlisting}[style=shadedpython]
NGRAM = 13
\end{lstlisting}
So it uses 13-word sequences.

\begin{paragraphblock}{Building banned 13-grams}
  Suppose an eval item contains:
\begin{lstlisting}[style=shadedpython,language={},basicstyle=\scriptsize\ttfamily]
the reserve bank of india was established in april nineteen thirty five in accordance
\end{lstlisting}
  A 13-word window might be hashed. Then the window shifts one word:
\begin{lstlisting}[style=shadedpython,language={},basicstyle=\scriptsize\ttfamily]
reserve bank of india was established in april nineteen thirty five in accordance with
\end{lstlisting}
  and so on. Each gets added to \texttt{banned}.
\end{paragraphblock}

\begin{paragraphblock}{Which eval fields are used?}
  For JSONL eval files it looks at \texttt{prompt}, \texttt{answer}, and
  \texttt{text} --- so both the question and answer content can contribute
  banned n-grams. For plain text eval files, the entire file contents are
  used.
\end{paragraphblock}

\begin{paragraphblock}{Strict contamination rule}

\begin{lstlisting}[style=shadedpython]
if any(g in banned for g in ngrams(doc["text"])):
\end{lstlisting}
  Meaning: if a training document shares even \emph{one} banned 13-word
  sequence with any eval content, drop the whole document.
\begin{lstlisting}[style=shadedpython,language={},basicstyle=\scriptsize\ttfamily]
training doc
   |
   v
generate all 13-grams
   |
   v
any overlap with eval?
  /           \\
 yes           no
  |             |
  v             v
drop           keep
\end{lstlisting}
  The rationale in the docstring: false positive / over-dropping is cheap;
  eval leakage is expensive --- especially at large scale.
\end{paragraphblock}

\subsubsection{Clean}

\begin{lstlisting}[style=shadedpython,language={},basicstyle=\scriptsize\ttfamily]
raw.jsonl
   |
   v
langid.py
   |
   | Adds:
   | * lang
   | * script
   | * code_mix
   |
   v
filters.py
   |
   | Rejects:
   | * too short / too long
   | * weird mean word length
   | * excessive symbols
   | * bullet-heavy text
   | * ellipsis-heavy text
   | * duplicate lines
   | * repeated bigrams
   | * too few alphabetic words
   | * language/stopword mismatch
   |
   v
dedup_minhash.py
   |
   | 5-word shingles
   |      |
   |      v
   | 128 MinHash values
   |      |
   |      v
   | 16 LSH bands
   |      |
   |      v
   | candidate duplicates
   |      |
   |      v
   | exact Jaccard >= 0.8
   |      |
   |      v
   | remove later duplicate
   |
   v
decontaminate.py
   |
   | build eval 13-grams
   |      |
   |      v
   | compare training docs
   |      |
   |      v
   | any shared 13-gram?
   |      |
   |      v
   | drop document
   |
   v
clean.jsonl
\end{lstlisting}

\begin{paragraphblock}{Code skips \texttt{filters.py}}
  This is especially clear from the actual filter rules. For example, code
  frequently contains \texttt{\#}, \texttt{*}, short lines, symbol-heavy
  expressions, few natural-language stopwords, unusual word lengths, and lots
  of nonalphabetic tokens. So a valid Python file like
\begin{lstlisting}[style=shadedpython]
def foo(x):
    if x > 10:
        return x * 2
\end{lstlisting}
  might fail several prose-oriented checks --- which is why code skips
  \texttt{filters.py} by design.

  Net result: 6{,}600 docs / $\sim$4.5M words, with the drops logged
  per stage (e.g.\ filters cut 402 of 3{,}300 Indic docs; decontamination
  banned 198 eval n-grams).
\end{paragraphblock}

\subsection{\texttt{eval\_fertility.py}}

For each \texttt{eval\_sets/<language>.txt}, the evaluator computes

\[
  \operatorname{fertility}
  =\frac{\text{number of encoded token IDs}}
  {\text{number of whitespace-delimited words}}.
\]

English becomes the relative reference. A language passes only if it is within
15\% of English; native-script languages must also remain at or below 1.8
tokens per word. Code is informational because it needs comparison with a
strong code tokenizer rather than the language threshold.

The current 64K result is:

\begin{center}
  \renewcommand{\arraystretch}{1.15}
  \small
  \begin{tabular}{@{}l r r l@{}}
    \textbf{Language} & \textbf{Fertility} & \textbf{vs English} &
    \textbf{Status} \\
    \hline
    English & 1.18 & 0\% & pass \\
    Hindi & 1.32 & +12\% & pass \\
    Punjabi & 1.67 & +42\% & fails relative target \\
    Bengali & 1.90 & +62\% & fails both targets \\
    Hinglish & 2.13 & +81\% & fails relative target \\
    Malayalam & 2.78 & +136\% & fails both targets
  \end{tabular}
\end{center}

The full evaluator reports 11 failures. These values diagnose the current
artifact, but the bundled texts are too small and are not valid held-out
evaluation sets. They cannot select the production tokenizer.

\subsection{What must be true before v0 freezes}

\begin{enumerate}
    \setlength{\itemsep}{0pt}
  \item Replace every placeholder evaluation file with versioned
    held-out text.
  \item Add at least one strong multilingual and one code-tokenizer baseline.
  \item Train the planned 64K/96K/128K candidates first; a one-month sprint
    should not jump directly to the old 160K/256K sweep without enough data.
  \item Compare downstream validation loss at equal model size and
    equal token budget.
  \item Add round-trip, special-token-ID, digit-split, and
    script-coverage tests.
  \item Freeze the tokenizer JSON, corpus manifest, evaluation hash,
    and decision record together.
\end{enumerate}

Fertility is necessary because it measures representation efficiency. It is
not sufficient because a tokenizer with excellent fertility can still produce
worse language-model loss or arithmetic behavior.
